\documentclass[../../main.tex]{subfiles}
\begin{document}
\section{Synthesis: what a proof of KLS now needs}
\label{sec:kls-synthesis}

The July 2026 preprints changed the shape of the problem, and it is worth stating the change
precisely rather than as a mood. \emph{The remaining difficulty is no longer quadratic forms or
third moments.} Conditional on \cite{Letwin2026QuadraticKLS}, $\kappa_n=O(1)$ and every quadratic
witness is eliminated. What remains is the problem of promoting fixed-matrix or averaged
information to uniform control of every nonlinear test function, without paying a logarithmic
largest-eigenvalue cost.

This section states that residue in six versions --- one per strategy family --- and then names the
four concrete next targets, each mapped to the open node of \texttt{research/kls/ledger.yaml} that
carries it.

\subsection{The residue, family by family}
\label{subsec:synthesis-table}

\begin{center}
\begin{tabularx}{\textwidth}{@{}>{\raggedright\arraybackslash}p{0.20\textwidth}>{\raggedright\arraybackslash}p{0.30\textwidth}Y@{}}
\toprule
Framework & What is now controlled & The missing estimate \\
\midrule
Needles (\S\ref{sec:family-needles}) & One-dimensional conditional measures, sharply &
A decomposition inheriting \emph{operator} covariance, not one or two scalar constraints \\
\addlinespace
Stochastic localization (\S\ref{sec:family-sl}) & Short-time covariance and, conditionally, sharp
third moments & Avoid the $\log n$ cost of approximating the maximum eigenvalue \\
\addlinespace
Heat flow, $H^{-1}$ (\S\ref{sec:family-bochner}) & Coordinate and quadratic spectral mass &
Uniform estimates for derivatives of an arbitrary $f$ \\
\addlinespace
Moment map (\S\ref{sec:family-moment-map}) & Fixed deterministic matrix energies
$\E\Tr(BHBH)$; and, exactly, $\CMH$ on the line, on products, and on every log-concave
Dirichlet law (\S\ref{sec:cmh-exact-cases}) & Control when the matrix or direction depends on
$X$ or on $f$. Even the linear sector is missing: $\lmax(\E H^2)\le4$ does not follow from
$\Tr(\E H^2)\le2n$, and Proposition~\ref{prop:letwin-not-gate-zero} shows no matrix-moment
argument closes the gap \\
\addlinespace
Parallel coupling (\S\ref{subsec:kls-solved-neighbours}) & Linear exponential tilts
$e^{\inner\theta x}\mu$ & Couplings for arbitrary functional perturbations \\
\addlinespace
Brownian transport (\S\ref{sec:family-transport}) & Polylogarithmic averaged derivative bounds &
A dimension-free expected operator derivative \\
\bottomrule
\end{tabularx}
\end{center}

Every row is the same sentence in a different dialect: something is controlled \emph{on average},
or \emph{for a fixed object}, and KLS needs it \emph{uniformly}, for an object that is allowed to
depend on the measure.

\subsection{The constraint any proposal must satisfy}
\label{subsec:synthesis-constraint}

Proposition~\ref{prop:covariance-spike} is not a technical nuisance; it is the sharpest available
guide to what a correct argument can look like. Products of centered exponentials have a
dimension-free Poincar\'e constant by tensorization, yet their conditional covariance spikes:
$\E\norm{\Cov(X\mid X+\sqrt sG)}_\op\gtrsim s$ for $s\le\log n$.

Two consequences follow, and they cut in opposite directions.

\begin{enumerate}[label=(\arabic*),leftmargin=2.4em]
  \item A straightforward ``bound $\norm{A_t}_\op$ better'' program cannot work: the statement it
  would need is false.
  \item Rare covariance spikes can be \emph{harmless}. A successful potential should recognize
  this rather than charge the full largest eigenvalue whenever a spike occurs. Tensorization is
  the diagnostic: a potential that is not tensorization-aware will charge $n$ independent
  coordinates $n$ times for a phenomenon that costs $O(1)$.
\end{enumerate}

\subsection{The four concrete next targets}
\label{subsec:synthesis-targets}

\paragraph{Target 1 --- function-adapted stochastic localization.}
For a fixed test function set $M_t(f)=\E_{p_t}f$, so that
\begin{equation}
  \dd M_t(f)=\Cov_{p_t}(X,f)\cdot\dd W_t .
  \label{eq:function-adapted-sde}
\end{equation}
Current proofs bound the integrand by the worst case,
\begin{equation}
  \abs{\Cov_{p_t}(X,f)}^2\le\norm{A_t}_\op\Var_{p_t}f ,
  \label{eq:worst-case-step}
\end{equation}
which discards essentially all information about $f$. A direct estimate of
\begin{equation}
  \int\frac{\abs{\Cov_{p_t}(X,f)}^2}{\Var_{p_t}f}\dd t
  \label{eq:function-adapted-occupation}
\end{equation}
for a near-extremizing eigenfunction would avoid the top-eigenvalue entropy cost of
\eqref{eq:logtraceexp} and would respect tensorization, since \eqref{eq:function-adapted-occupation}
factorizes over independent blocks in a way that $\norm{A_t}_\op$ does not.

\emph{Repository target:} this is exactly Question~\ref{q:mm-spectral-occupation}, the headline of
Route~S; the fixed-cut analogue is \texttt{q:upgrade}, the operator-to-trace upgrade of
Section~\ref{sec:open}.

\paragraph{Target 2 --- a nonlinear extension of the moment-map estimate.}
Brascamp--Lieb in moment-map coordinates already gives \eqref{eq:mm-brascamp-lieb}, so KLS would
follow from
\begin{equation}
  \E\inner{\tau_\mu\nabla f}{\nabla f}\lesssim\E\abs{\nabla f}^2 .
  \label{eq:nonlinear-moment-map}
\end{equation}
The identity $\E\tau_\mu=I$ is insufficient, because $\tau_\mu(X)$ may correlate with
$\nabla f(X)$. Letwin controls a deterministic $B$; the missing theorem must handle an
$X$-dependent direction or matrix field.

\emph{Repository target:} Route~C, Group~5 --- \texttt{prog:cmh-route}. Its regular-class
endpoint is now normalized: Definition~\ref{def:cmh} fixes the operator data and
Theorem~\ref{thm:cmh-implies-affine-poincare} proves $\CPaff\le\CMH$, discharging
\texttt{q:cmh-normalization}; passage to arbitrary log-concave limits is the separate open node
\texttt{q:cmh-approximation}. The route's other open sub-targets split into a construction layer
(\texttt{q:mm-invariant-lift}, \texttt{q:mm-square-root-commutator}) and a falsification layer
(\texttt{q:gate-zero}, \texttt{q:cmh-solenoidal-perturbation}).

Two cautions belong in this table rather than in the route brief, because they change how this
target should be read. First, $\mathrm{CMH}(4)$ is \emph{not} a reformulation of
\eqref{eq:nonlinear-moment-map} or of Conjecture~\ref{conj:kls}: by
Proposition~\ref{prop:cmh-hodge} it additionally demands control of a solenoidal excess that
vanishes identically in dimension one under the no-flux convention. This proves an additional
channel, but not a strict-separation theorem; equivalence at constant $4$ remains open
(Remark~\ref{rem:cmh-stronger-than-kls}). Second, its cheapest necessary consequence, gate zero
(Conjecture~\ref{conj:gate-zero}), is itself an instance of the row above: in isotropic position
it asks $\lmax(\E H^2)\le4$ where Theorem~\ref{thm:chen-klartag-moment-hessian} supplies only the
trace bound $\Tr(\E H^2)\le2n$. Target~2 therefore does not escape the average-versus-uniform
pattern; it relocates it, and Proposition~\ref{prop:letwin-not-gate-zero} shows the relocation is
not free.

\paragraph{Target 3 --- replace log-trace-exp by an effective-rank estimate.}
By Remark~\ref{rem:log-is-entropy}, the remaining $\log n$ enters \emph{solely} because the soft
maximum \eqref{eq:logtraceexp} approximates $\lmax$ over $n$ directions. A potential depending
only on the directions actually relevant to a near-extremizer --- or on an effective rank rather
than the ambient dimension --- would convert $\kappa_n=O(1)$ directly into $\CP=O(1)$.

\emph{Repository target:} the interface functional $\Xi_T(\mu)$ of \eqref{eq:interface-def} and
its evaluation in Section~\ref{sec:bootstrap} are this repository's version of the question, and
\texttt{q:taming} is the corresponding open node. Section~\ref{sec:bootstrap} already shows that
the crude evaluation is provably insufficient and that a relative bound at a sufficiently small
universal time is KLS-sufficient.

\paragraph{Target 4 --- extend parallel coupling beyond linear tilts.}
Present parallel coupling controls the finite-dimensional family $e^{\inner\theta x}\mu(\dd x)$. A
coupling for perturbations $(1+\eps f)\mu$ with cost controlled by $\int\abs{\nabla f}^2\dd\mu$
would address arbitrary spectral directions directly, rather than one linear family.

\emph{Repository target:} none currently. This is a gap in the route registry, not a gap in the
literature survey, and it is recorded here as such.

\subsection{Assessment}
\label{subsec:synthesis-assessment}

Of the four, target~1 is the best aligned with the known obstruction: it is the only one that
avoids asking for a uniform top-covariance statement, and by
Section~\ref{subsec:synthesis-constraint} such a statement is false for tensorized exponentials.
That is a structural argument for its priority, not merely a preference.

Target~2 is conceptually deeper and may ultimately be cleaner, since
\eqref{eq:nonlinear-moment-map} is a single inequality with no stochastic apparatus at all. But it
requires controlling correlations with nonlinear gradient fields, which is a considerably stronger
theorem than Letwin's fixed-matrix estimate --- and the repository's own experience on Route~C
(Section~\ref{sec:moment-map-cmh}) is that the difficulty concentrates in a commutator sum that
has so far resisted dimension-free control.

Target~3 is the most narrowly technical and the most clearly delimited: the quantity to be removed
is identified exactly, in one displayed inequality. It is also the one for which this repository
has already proved the most --- including the negative result that the crude evaluation cannot
suffice.

A final caution on status. Three of the four targets take $\kappa_n=O(1)$ as their starting point,
and that input is an unreviewed version-1 preprint (Section~\ref{subsec:mm-audit}). Should it not
survive review, targets~1--3 do not become wrong, but their premise reverts to
$\CP\lesssim\log n$ and the arithmetic of every ``remaining gap'' claim in this section changes.
The ledger records this dependency explicitly through \texttt{import\_class:
preprint-unreviewed}, which cannot underlie a node marked \texttt{proved}; that mechanism, and not
this prose, is what keeps the distinction enforced.
\end{document}
